\chapter{Stage VII：exceptional point と幾何学的 holonomy}

\solproblem{VII.1　2-by-2 complex-symmetric EP}{
一般 complex-symmetric 行列を対角化し二次 EP 条件を求めよ。}{
\begin{equation}
 H=\begin{pmatrix}a&b\\b&d\end{pmatrix},\qquad
 E_\pm=\frac{a+d}{2}\pm\frac12\sqrt{(a-d)^2+4b^2}.
\end{equation}
従って discriminant
\begin{equation}
 \Delta=(a-d)^2+4b^2=0
\end{equation}
で固有値が合流する。さらに $b=0$ かつ $a=d$ なら $H=a\identity$ で二本の
固有vector が残る ordinary degeneracy である。EP には
$H-\vsub{E}{EP}\identity\neq0$、rank 1 が必要。例えば
$r=(b,E-a)^{\mathsf T}$ の二枝は $\Delta\to0$ で同一直線へ合流する。}

\solproblem{VII.2　Jordan chain と resolvent の二重 pole}{
長さ2の Jordan chain を作り resolvent の $x^{-2},x^{-1}$ 項を求めよ。}{
$(H-z_0)r_0=0$、$(H-z_0)r_1=r_0$ とし、Jordan block 上で
$H=z_0\identity+N$、$N^2=0$。$x=z-z_0$ なら本書の resolvent は
\begin{equation}
 (z-H)^{-1}=(x\identity-N)^{-1}
 =\frac{\identity}{x}+\frac{N}{x^2}.
\end{equation}
dual Jordan basis $l_0,l_1$ を
$\dualpair{l_i}{r_j}=\delta_{ij}$ と取れば
$\identity_J=\ket{r_0}\bra{l_0}+\ket{r_1}\bra{l_1}$、
$N=\ket{r_0}\bra{l_1}$。従って
\begin{equation}
 \Resolvent_J(z)=\frac{\ket{r_0}\bra{l_1}}{x^2}
 +\frac{\ket{r_0}\bra{l_0}+\ket{r_1}\bra{l_1}}{x}.
\end{equation}
$\Cincoming(E)$ の二重零点の逆に現れる二係数と同じ pole order である。}

\solproblem{VII.3　EP の self-orthogonality}{
coalesced 右 vector と左 partner の内積が0になることを示せ。}{
EP で右 Jordan chain を
$(H-E_*)r_0=0$、$(H-E_*)r_1=r_0$、左固有vectorを
$l_0(H-E_*)=0$ とする。第二式を左から掛ければ直ちに
\begin{equation}
 \dualpair{l_0}{r_0}=\dualpair{l_0}{(H-E_*)r_1}=0.
\end{equation}
parameter 微分でも同じ結果が見える。
$(H-E)r=0$ を $\lambda$ で微分して左固有vectorを掛けると
$E'\dualpair l r=\dualpair l{H'r}$。EP 近傍で
$E'\sim(\lambda-\lambda_*)^{-1/2}$ が発散する一方、右辺が有限なら
$\dualpair l r\to0$ が必要である。}

\solproblem{VII.4　branch point の一周と二周}{
square-root EP を周回して固有値・固有vector・gauge を追跡せよ。}{
局所 normal form $H(\lambda)=\begin{pmatrix}0&1\\\lambda&0\end{pmatrix}$ を
用いる。$\lambda=\rho\rme^{\rmi\varphi}$ とすると
$E_\pm=\pm\sqrt\rho\rme^{\rmi\varphi/2}$、
$r_\pm=(1,E_\pm)^{\mathsf T}$。$\varphi:0\to2\pi$ で
$E_+\to E_-$、$r_+\to r_-$、すなわち label が交換する。二周で固有値は
元へ戻る。vector の比較には各 step で biorthogonal overlap を正実にする
parallel-transport gauge を選び、最後だけ初期 gauge と比較する。任意に
各点の第一成分を実にする gauge は途中で不連続となり位相を捏造する。
二周後の符号または位相は、この連続 gauge で初めて幾何学的 holonomy と
して定義される。}

\solproblem{VII.5　biorthogonal Berry connection}{
二準位 model の Berry connection と gauge 変換部分を求めよ。}{
上の model で右 vector $r=(1,E)^{\mathsf T}$、左 vector を
$l^{\mathsf T}=(E,1)/(2E)$ とすれば $l^{\mathsf T}r=1$。従って
\begin{equation}
 \mathcal A=\rmi l^{\mathsf T}\rmd r
 =\frac{\rmi}{2}\frac{\rmd E}{E}
 =\frac{\rmi}{4}\frac{\rmd\lambda}{\lambda}.
\end{equation}
$r\mapsto\rme^{f}r$、$l\mapsto\rme^{-f}l$ では
$\mathcal A\mapsto\mathcal A+
\rmi\rmd f$。従って局所 connection は gauge
依存だが、閉路積分の exponential は $f$ が一価なら不変。EP の一周では
state 自体が相手の sheet へ移るので、二周または permutation matrix を
含む non-Abelian transport として比較する。}

\solproblem{VII.6　momentum bin と continuum normalization}{
delta-normalized state を幅 $\Delta k$ の bin に置き換え、発散 normalization と有限 holonomy を分離せよ。}{
$\braket{k|k'}=\delta(k-k')$ とし
\begin{equation}
 \ket{n}=\frac1{\sqrt{\Delta k}}
 \int_{k_0-\Delta k/2}^{k_0+\Delta k/2}\ket{k}\dd k.
\end{equation}
なら $\braket{n|n}=1$。一方、中央の generalized state に戻す係数は
$\ket{n}\simeq\sqrt{\Delta k}\ket{k_0}$、従って
$\ket{k_0}\simeq(\Delta k)^{-1/2}\ket{n}$ と発散する。parameter-dependent
bin の Berry connection は
\begin{equation}
 \dualpair{\widetilde n}{\rmd n}
 =-\frac12\rmd\log\Delta k+\vsub{\mathcal A}{shape}+\vsub{\mathcal A}{phys}.
\end{equation}
最初の exact differential は normalization artifact。閉路で bin の幅と
端点を同じ値へ戻し、shape 項を同一規約で差し引けば
$\oint\mathcal \vsub{A}{phys}$ が有限 holonomy として残る。}

\solproblem{VII.7　bin-width independence の RG 方程式}{
renormalized geometric phase の $\Delta k$ 非依存条件と積分定数を求めよ。}{
bare phase を $\gamma_B(\Delta k)$、counterterm を $C(\Delta k,\mu)$ とし
\begin{equation}
 \gamma_R(\mu)=\gamma_B(\Delta k)+C(\Delta k,\mu),\qquad
 \Delta k\frac{\rmd\gamma_R}{\rmd\Delta k}=0
\end{equation}
を課す。もし $\Delta k\rmd\gamma_B/\rmd\Delta k=c+o(1)$ なら
\begin{equation}
 C=-c\log\frac{\Delta k}{\mu}+C_0,\qquad
 \mu\frac{\rmd\gamma_R}{\rmd\mu}=c.
\end{equation}
$C_0$ は局所発散からは決まらない。reference loop の phase を0とする、
detector の asymptotic phase convention を固定する、または一つの実測値へ
match する、といった物理的 convention/data がこれを決める。}

\solproblem{VII.8　EP と continuum threshold の競合}{
EP が threshold へ近づくとき二つの branch をどう扱うべきか説明せよ。}{
threshold $\vsub{E}{th}$ では resolvent 自己 energy が
$\Sigma(E)\sim a+b(E-\vsub{E}{th})^{l+1/2}$ を持つ。一方孤立 EP は
$E_\pm-E_*\sim\pm c\sqrt{\lambda-\lambda_*}$。両者が近いと局所 determinant
は例えば
\begin{equation}
 D(E,\lambda)=A(E-E_*)^2+B(\lambda-\lambda_*)
 +C(E-\vsub{E}{th})^{l+1/2}+\cdots
\end{equation}
となり、単純な二次多項式 normal form は一様でない。$s$ wave なら
$k=\sqrt{E-\vsub{E}{th}}$ で uniformize し、$D(k,\lambda)=0$ と
$\partial_kD=0$ を同じ sheet で解く。EP sheet exchange と $k\mapsto-k$ の
threshold exchange が合成され、周回 permutation は経路がどの branch
point を囲むかに依存する。}

\solproblem{VII.9　condition number と pseudospectrum}{
EP 周回上の eigenvalue condition number を計算し pseudospectral sensitivity と結べ。}{
単純固有値の condition number は
\begin{equation}
 \kappa_n=\frac{\|r_n\|\|l_n\|}{|\dualpair{l_n}{r_n}|}.
\end{equation}
biorthonormal gauge でも vector norm の積として同じ値が残る。二次 EP から
距離 $\rho=|\lambda-\lambda_*|$ では左右 overlap が
$\order(\sqrt\rho)$ なので $\kappa\sim\rho^{-1/2}$。摂動 $\delta H$ に対し
$|\delta E_n|\lesssim\kappa_n\|\delta H\|$。従って周回上で $\kappa$ が大きい
点は $\epsilon$-pseudospectrum の等高線が膨らむ点と一致する。数値的には
左右 residual を一定に保って $\kappa(\varphi)$ と最小 singular value
$s_{\min}(E-H)$ を同時に plot する。}

\solproblem{VII.10　multiple-zero continuation test}{
$\mathcal C=\partial_E\mathcal C=0$ を用いる探索と near crossing との識別法を設計せよ。}{
scalar determinant または適切な incoming minor について
\begin{equation}
 F_1(E,\lambda)=\mathcal C(E,\lambda)=0,\qquad
 F_2(E,\lambda)=\partial_E\mathcal C(E,\lambda)=0
\end{equation}
を complex Newton 法で同時に解く。Jacobian は
$\partial(E,\lambda)(F_1,F_2)$、必要なら complex-step/automatic
differentiationを使う。候補点の検証は
\begin{enumerate}[label=(\alph*)]
 \item operator の geometric multiplicity が1、Jordan residual が小さい；
 \item 小円周回で二 pole が交換し、二周で戻る；
 \item splitting が radius の平方根に比例；
 \item basis・contour・complex-scaling angle を変えて点と Riesz projector が安定；
 \item $\mathcal C_{EE}\neq0$、$\mathcal C_\lambda\neq0$。
\end{enumerate}
near crossing は有限最小 gap を持ち、一周で label 交換せず、Newton 解が
discretization とともに漂う。}
